\chapter{Considerações Finais}\label{chp:CONCLUSAO}

Este trabalho teve como objetivo o desenvolvimento de um sistema \textit{web} para a reserva de mesas, salas e auditórios da universidade, dotado de uma interface de visualização baseada nas plantas baixas reais dos ambientes. A partir do levantamento das necessidades da comunidade acadêmica e do estudo das tecnologias pertinentes, foi projetada e implementada uma solução funcional, organizada em uma arquitetura em camadas e apoiada em tecnologias web consolidadas.

Os resultados apresentados no Capítulo~\ref{chp:RESULTADOS} indicam que o sistema atingiu os objetivos propostos: a reserva de recursos pode ser realizada de forma autônoma e visual, a disponibilidade é verificada de modo a evitar conflitos, e o fluxo de aprovação, associado à política de penalidades por não comparecimento, promove o uso responsável dos espaços. A representação fiel dos ambientes, viabilizada pelas plantas baixas cedidas pela COPEA e pelo processo de vetorização e normalização dos desenhos, constitui o principal diferencial da solução, aproximando a experiência digital do ambiente físico.

Como contribuição, o trabalho oferece uma alternativa concreta para a gestão de espaços acadêmicos, com potencial de reduzir conflitos, otimizar a ocupação da infraestrutura e aumentar a transparência do processo de reserva. Confirma-se, assim, a hipótese de que um sistema de reservas com interface de visualização realista favorece uma gestão mais eficiente dos espaços da universidade.

\section{Trabalhos Futuros}\label{sec:trab_futuros}

A partir dos resultados obtidos e das limitações identificadas, apontam-se as seguintes possibilidades de trabalhos futuros:

\begin{itemize}
  \item Ampliar a validação com usuários, por meio de testes de usabilidade em maior escala, a fim de consolidar as conclusões sobre a experiência de uso;
  \item Estender a cobertura do sistema a outros pavilhões e edifícios da universidade, incorporando novas plantas baixas;
  \item Automatizar a etapa de normalização dos vetores, reduzindo o esforço manual na preparação de novas plantas;
  \item Desenvolver uma versão para dispositivos móveis ou um aplicativo dedicado, ampliando a acessibilidade da solução;
  \item Incorporar recursos de análise de dados e relatórios gerenciais mais elaborados, apoiando a tomada de decisão sobre a utilização dos espaços;
  \item Integrar o sistema a serviços institucionais de autenticação, unificando o acesso com as credenciais já utilizadas pela comunidade acadêmica.
\end{itemize}

\pagebreak